%% Preamble — A First Course in Monte Carlo Methods
%% D. Sanz-Alonso & O. Al-Ghattas, University of Chicago

\documentclass[12pt,openany]{book}

\usepackage[T1]{fontenc}
\usepackage[utf8]{inputenc}
\usepackage{lmodern}
\usepackage{microtype}

%% Page geometry — matches source
\usepackage[
  paperwidth=6.5in,
  paperheight=9.5in,
  top=1in,
  bottom=1in,
  inner=1in,
  outer=0.85in,
  headheight=14pt,
  headsep=0.3in
]{geometry}

%% Math
\usepackage{amsmath}
\usepackage{amssymb}
\usepackage{amsthm}
\usepackage{mathtools}
\usepackage{bm}
\usepackage{bbm}        % \mathbbm{1} for indicator function

%% Graphics
\usepackage{graphicx}
\graphicspath{{../figures/}{./figures/}}

%% TikZ (for Markov chain diagrams)
\usepackage{tikz}
\usetikzlibrary{arrows.meta,positioning,automata}

%% Algorithms
\usepackage{algorithm}
\usepackage{algpseudocode}
\renewcommand{\algorithmicrequire}{\textbf{Input:}}
\renewcommand{\algorithmicensure}{\textbf{Output:}}
\algnewcommand{\LineComment}[1]{\State \(\triangleright\) #1}

%% Lists and tables
\usepackage[shortlabels]{enumitem}
\usepackage{makeidx}
\makeindex
\usepackage{booktabs}
\usepackage{array}

%% Colors and links
\usepackage{xcolor}
\usepackage[
  colorlinks=true,
  linkcolor={blue!60!black},
  citecolor={blue!60!black},
  urlcolor={blue!60!black}
]{hyperref}

%% Section heading style: "1.1 · Section Title"
\usepackage{titlesec}
\titleformat{\section}[block]
  {\large\bfseries}
  {\thesection\ $\cdot$\ }
  {0pt}{}
\titleformat{\subsection}[block]
  {\normalsize\bfseries}
  {\thesubsection\ $\cdot$\ }
  {0pt}{}

%% Suppress page numbers on blank pages inserted by \cleardoublepage
\makeatletter
\renewcommand\cleardoublepage{\clearpage\if@twoside\ifodd\c@page\else
  \hbox{}\thispagestyle{empty}\newpage\if@twocolumn\hbox{}\newpage\fi\fi\fi}
\makeatother

%% Headers/footers
\usepackage{fancyhdr}
\pagestyle{fancy}
\fancyhf{}
\fancyhead[LE]{\small\thepage\quad\nouppercase{\leftmark}}
\fancyhead[RO]{\small\nouppercase{\rightmark}\quad\thepage}
\renewcommand{\headrulewidth}{0.4pt}

%% Spacing
\setlength{\parindent}{1.4em}
\setlength{\parskip}{0pt}
\linespread{1.05}
\raggedbottom
\emergencystretch=2em

%% Equation / figure / algorithm numbering by chapter
\numberwithin{equation}{chapter}
\numberwithin{figure}{chapter}
\numberwithin{table}{chapter}
\numberwithin{algorithm}{chapter}

%% -------------------------------------------------------
%% Theorem environments — shared counter per chapter
%% -------------------------------------------------------
\theoremstyle{plain}
\newtheorem{theorem}{Theorem}[chapter]
\newtheorem{proposition}[theorem]{Proposition}
\newtheorem{lemma}[theorem]{Lemma}
\newtheorem{corollary}[theorem]{Corollary}

\theoremstyle{definition}
\newtheorem{definition}[theorem]{Definition}
\newtheorem{example}[theorem]{Example}
\newtheorem{assumption}[theorem]{Assumption}

\theoremstyle{remark}
\newtheorem{remark}[theorem]{Remark}

%% -------------------------------------------------------
%% Custom math commands
%% -------------------------------------------------------

% Probability, expectation, variance
\newcommand{\Prob}{\mathbb{P}}
\newcommand{\Expect}{\mathbb{E}}
\newcommand{\Var}{\operatorname{Var}}
\newcommand{\Cov}{\operatorname{Cov}}

% Number fields
\newcommand{\R}{\mathbb{R}}
\newcommand{\N}{\mathbb{N}}
\newcommand{\Z}{\mathbb{Z}}

% Distributions
\newcommand{\Normal}{\mathcal{N}}
\newcommand{\Unif}{\operatorname{Unif}}
\newcommand{\Exp}{\operatorname{Exponential}}
\newcommand{\Beta}{\operatorname{Beta}}
\newcommand{\Binom}{\operatorname{Binomial}}
\newcommand{\Geom}{\operatorname{Geometric}}
\newcommand{\GammaDist}{\operatorname{Gamma}}
\newcommand{\Bern}{\operatorname{Bernoulli}}

% Indicator function
\newcommand{\Ind}[1]{\mathbbm{1}_{\{#1\}}}
\newcommand{\ind}{\mathbbm{1}}

% Integral functional notation from the book: \calI_f[h]
\newcommand{\calI}{\mathcal{I}}
\newcommand{\calF}{\mathcal{F}}
\newcommand{\calG}{\mathcal{G}}
\newcommand{\calB}{\mathcal{B}}
\newcommand{\calX}{\mathcal{X}}
\newcommand{\calY}{\mathcal{Y}}
\newcommand{\calZ}{\mathcal{Z}}
\newcommand{\calA}{\mathcal{A}}
\newcommand{\calL}{\mathcal{L}}
\newcommand{\calO}{\mathcal{O}}
\newcommand{\calP}{\mathcal{P}}
\newcommand{\calN}{\mathcal{N}}

% Bold vectors / matrices
\newcommand{\bfx}{\mathbf{x}}
\newcommand{\bfy}{\mathbf{y}}
\newcommand{\bfz}{\mathbf{z}}
\newcommand{\bfq}{\mathbf{q}}
\newcommand{\bfp}{\mathbf{p}}
\newcommand{\bftheta}{\boldsymbol{\theta}}
\newcommand{\bfmu}{\boldsymbol{\mu}}
\newcommand{\bfSigma}{\boldsymbol{\Sigma}}
\newcommand{\bfphi}{\boldsymbol{\phi}}

% KL divergence and variational inference
\newcommand{\KL}[2]{D_{\mathrm{KL}}\!\left(#1\,\|\,#2\right)}
\newcommand{\ELBO}{\mathrm{ELBO}}

% Norms
\newcommand{\abs}[1]{\left|#1\right|}
\newcommand{\norm}[1]{\left\|#1\right\|}

% Operators
\DeclareMathOperator*{\argmax}{arg\,max}
\DeclareMathOperator*{\argmin}{arg\,min}
\newcommand{\iid}{\overset{\mathrm{iid}}{\sim}}
\newcommand{\cvd}{\xrightarrow{d}}
\newcommand{\cvp}{\xrightarrow{p}}
\newcommand{\cvas}{\xrightarrow{\mathrm{a.s.}}}

% Gradient / Laplacian
\newcommand{\grad}{\nabla}
\newcommand{\Laplacian}{\Delta}

% Hamiltonian dynamics notation
\newcommand{\Ham}{\mathcal{H}}

% Acceptance probability
\newcommand{\accprob}{\alpha}

%% TOC depth: sections and subsections (matches source TOC)
\setcounter{tocdepth}{2}
